\section{Final-Checkpoint Performance}
\label{sec:results_final_checkpoint}

Final-checkpoint return measures the quality of the policy at the end of the fixed training budget.
It complements step-AUC by emphasizing asymptotic performance within the allotted interactions.
Table~\ref{tab:results_final_checkpoint} reports final deterministic evaluation returns for GLPE, GLPE without gating, and the strongest baseline in each environment.

\begin{table}[!htbp]
\centering
\small
\begin{fitwidthtabular}{lccc}
\hline
Environment & GLPE & GLPE (no gate) & Best baseline \\
\hline
\textsc{MountainCar-v0} & $-96.2\;[-96.5,-95.9]$ & $-97.1\;[-99.0,-96.0]$ & RIAC: $-106.6\;[-127.5,-95.9]$ \\
\textsc{BipedalWalker-v3} & $266.5\;[214.0,301.3]$ & $294.4\;[285.9,303.7]$ & PPO: $302.2\;[291.2,309.4]$ \\
\textsc{Ant-v5} & $4{,}961\;[4{,}421,5{,}359]$ & $4{,}961\;[4{,}436,5{,}349]$ & RIAC: $5{,}384\;[5{,}118,5{,}625]$ \\
\textsc{HalfCheetah-v5} & $6{,}889\;[6{,}565,7{,}285]$ & $6{,}908\;[6{,}548,7{,}280]$ & RIDE: $7{,}106\;[6{,}901,7{,}299]$ \\
\textsc{Humanoid-v5} & $1{,}970\;[546,4{,}355]$ & $2{,}858\;[1{,}428,4{,}454]$ & ICM: $2{,}874\;[770,4{,}956]$ \\
\hline
\end{fitwidthtabular}
\caption{Final-checkpoint mean extrinsic return under deterministic offline evaluation. Brackets show 95\% bootstrap confidence intervals over seeds. Higher values are better.}
\label{tab:results_final_checkpoint}
\end{table}

The final-checkpoint results reinforce the environment-dependent character of the method.
GLPE is the strongest method on \textsc{MountainCar-v0}, where the return interval is narrow and close to the maximum observed solution quality.
The reliability of this result is important because \textsc{MountainCar-v0} contains limited extrinsic feedback until the agent has already learned an effective control strategy.

In \textsc{BipedalWalker-v3}, the ungated variant is much closer to the strongest baseline than gated GLPE.
The difference suggests that gate suppression can remove useful intrinsic shaping when the environment's dynamics are learnable but the progress signal fluctuates near the threshold.
The same pattern appears more strongly in \textsc{Humanoid-v5}, where GLPE without gating is close to ICM while gated GLPE is lower.

In \textsc{Ant-v5} and \textsc{HalfCheetah-v5}, the GLPE family achieves high returns but does not exceed the strongest baseline.
These tasks provide dense reward components, so exploration bonuses mainly influence the path of optimization rather than enabling reward discovery from near-zero feedback.
The results therefore suggest that GLPE is most appropriate when the exploration signal must compensate for sparse or delayed extrinsic feedback.